\section{Implementation and Integration Status}
\label{sec:implementation}

\ifanonymouspreview
The source audit fixes separate file-system and object-backend research
revisions; the internal evidence ledger retains their identities and source
locations.
\else
The source audit fixes \sysname{} at \code{242af5f7e9bb} and the RustFS
research fork at \code{cc4926fac110}.\footnote{Fixed sources:
\url{https://github.com/brewfs/brewfs/tree/242af5f7e9bba7afc0c61778f1084fcb20e083a0} and
\url{https://github.com/Ivanbeethoven/rustfs/tree/cc4926fac1107aee76540f28bf22dbd47fcbca24}.}
\fi
These revisions define the implementation scope below. Source inspection
establishes which components and connections exist; it does not establish
successful deployment or measured system behavior. This revision adds no
mount, fault-injection, or performance experiments.

\paragraph{Mount entry points.}
Table~\ref{tab:implementation-entrypoints} separates the current paths.
Their separation realizes optional mechanisms without implying that every
combination is available through every entry point.

\begin{table*}[t]
\centering
\small
\caption{Current mount entry points and their construction. Compile-time
features and backend capabilities further restrict admission.}
\label{tab:implementation-entrypoints}
\begin{tabularx}{\textwidth}{@{}>{\raggedright\arraybackslash}p{0.25\textwidth}Y@{}}
\toprule
Entry point & Constructed file-system view \\
\midrule
Ordinary \code{flat-v1} &
Ordinary metadata store and VFS with an object-data path; no workspace
upper/lower construction. \\
\code{workspace-v1} &
One sealed catalog base and one private writable upper; an optional
persisted binding selects a Packed v3 lower adapter. \\
Standalone Packed v3 &
Authenticated packed snapshot and read-only VFS; no private writable
upper is required. \\
Native P1 (\code{workspace-native-v2}) &
Workspace-v1 KV namespace and base-data source, with Native data runtime
and ordinary VFS. \\
Experimental Entity (\code{workspace-entity-v1}) &
Fresh Redis-backed current-entity namespace over a packed manifest,
S3 data, and a dedicated FUSE adapter. \\
\bottomrule
\end{tabularx}
\end{table*}

The \code{workspace-v1} catalog entry accepts SQLite, Redis, and TiKV, but
rejects etcd. Attaching a Packed lower additionally requires the catalog's
permission, identifier, and lifecycle capabilities. The default
\code{WorkspaceLower::Native} means native catalog layers; it is distinct
from the \code{native\_base} subsystem. Native P1 requires a workspace-v1
schema-1 base and a KV namespace. Its CLI explicitly rejects a frozen
metadata volume rather than treating the mutable KV path as a fallback.
Library support for frozen Native metadata therefore does not establish
that this mount entry supports it.

\paragraph{Reader and Entity bounds.}
The standalone Packed entry accepts the current wire005 format through a
content-addressed manifest; the CLI rejects earlier Packed v1/v2 entries.
It also rejects nonzero payload memory-cache, payload SSD-cache, and
prefetch budgets. Supported metadata caches must be distinguished from
these unavailable payload-cache configurations.

Entity is an experimental, fresh-only path with process-crash takeover
disabled. Its adapter uses direct I/O, zero metadata TTL, and no ordinary
VFS slice-cache or writeback path. The Entity data API bounds an individual
read, write, and logical chunk size at 16 MiB; the FUSE adapter separately
limits writes to 1 MiB. Its complete-directory resolver bounds enumeration
at 4,096 entries. Durable Entity lower ownership is limited to 256 active
owner records per catalog store, with admission and census checks. This is
distinct from the separate bounded set of Packed reader leases. These are
path-specific constraints, not global file-size, directory-size, or workspace
limits for all mounts. Entity upper reads currently fetch and authenticate
the complete stored object before returning the requested subrange.

\paragraph{Publication and recovery.}
The code includes authenticated snapshot production, persisted lower
bindings, reader pins, and guarded publication. It also wires a restricted
Packed recovery command for reusing the same persistent volume claim
(PVC). That path checks recovery generation, lease and pod identities, and
scoped metadata/PVC authority. It does not provide Entity takeover or
automatic migration to an arbitrary worker.

Native seal, publish, and fork functions encode quiescence, object
verification, retention, and atomic revision transitions. The inspected
tree contains these lifecycle definitions and unit-test callers, but no
ordinary CLI handoff driver. Consequently, their protocol structure
does not demonstrate an integrated worker transfer. Recovering a published
file state also does not recover process memory or open application
sessions; applications must establish suitable checkpoints and reopen
their tools.

\paragraph{Object formation and layout gates.}
The ordinary object path fixes client block size $g$ in configuration,
with a default of 4 MiB inside 64 MiB logical chunks. Readers use the
declared geometry, and incompatible volume configuration is rejected.
There is no implemented access-driven controller selecting the proposed
1/4/8/16 MiB client sizes. Changing $g$ on existing data requires preserving
its interpretation, not simply changing a mount option.

The encoding-block-size hint requires an explicit compile-time feature,
is off by default, and requires a compatible S3 backend. It does not change
$g$ or multipart part size. The inspected backend independently gates
acceptance of $B$ hints and creation of new $q$ framing, each with enable
and fleet-confirmation gates that default to off. Eligibility and geometry
checks still apply. Reads and repair use the object's persisted layout;
enabling a client hint does not re-encode existing objects.

\paragraph{Remaining connections and experimental controls.}
The ordinary uploader propagates configured hints, but the Entity mount
constructs its S3 client with the default hint setting, and its strict
durable create-only upload does not attach hint metadata. The Packed
fixture producer also uses the default setting. Extending these paths must
preserve conditional creation and precondition-failure behavior. Moreover,
Entity directly uploads segments of each FUSE write rather than aggregating
ordinary client blocks. Its upper objects therefore remain at most 1 MiB
on this path and do not trigger the automatic hint policy.

That automatic hint policy omits a hint for $L_o\leq1\MiB$ and requests
$B=1\MiB$ otherwise. The inspected research backend already defaults to
$B=1\MiB$, so enabling the policy can leave effective $B$ unchanged.
Hint configuration also changes the ordinary client's fetch and cache
policy, even for objects receiving no hint. Experiments must therefore
record actual $(g,L_o,B,q)$, eligibility, cache behavior, and physical
read bytes, and isolate these factors. A configuration named \emph{all-on}
alone proves neither a complete combined runtime nor a layout treatment.
